\lysh{16773_SOLUTION}{1}
\hfill \break
ΛΥΣΗ
\hfill \break
α) Η εξίσωση κύκλου με κέντρο το $Ο(0,0)$ είναι $x^2 + y^2 = \rho^2$ (1), όπου $\rho$ η ακτίνα του κύκλου.

Επειδή ο κύκλος διέρχεται από το σημείο $A(1,2)$ θα πρέπει οι συντεταγμένες του $A$ να επαληθεύουν την εξίσωση (1), δηλαδή $1^2 + 2^2 = \rho^2$, άρα $\rho^2 = 5$.

Οπότε, η (1) γίνεται:
\[ x^2 + y^2 = 5 \]

β)
i. Η εξίσωση της εφαπτομένης του κύκλου (1) στο σημείο του $A(x_1, y_1)$ είναι:
\[ xx_1 + yy_1 = \rho^2 \quad (2) \]
Έτσι η εξίσωση της εφαπτομένης του κύκλου $x^2 + y^2 = 5$ στο σημείο του $A(1,2)$ θα είναι:
\[ x + 2y = 5 \]

ii. Για να είναι το σημείο $B$ αντιδιαμετρικό του $A$, θα πρέπει το κέντρο $O$ να είναι το μέσο του τμήματος $AB$. Επομένως θα ισχύουν:
\[
\left\{
\begin{array}{l}
x_O = \dfrac{x_A + x_B}{2} \\
y_O = \dfrac{y_A + y_B}{2}
\end{array}
\right.
\text{ ή }
\left\{
\begin{array}{l}
0 = \dfrac{1 + x_B}{2} \\
0 = \dfrac{2 + y_B}{2}
\end{array}
\right.
\text{ ή }
\left\{
\begin{array}{l}
0 = 1 + x_B \\
0 = 2 + y_B
\end{array}
\right.
, \text{ άρα } x_B = -1 \text{ και } y_B = -2.
\]
Άρα είναι $B(-1,-2)$.
\hfill \break
\vspace{10pt}
\hfill \break